\chapter{Transient Simulations}
\label{ch:transient}

% local: velocity unknown in the first-order recast (no \vv collision; \vv is defined in simbook)

\chapterorient{We return to the anchor of \cref{ch:electrostatics} one last time,
and this time we swap the corner that the rest of Part~V has been circling. The
electrostatic and driven analyses solve a \emph{family} of problems---a fan of
independent solves, a \texttt{map}, parallel over terminals or frequencies. The
transient analysis does something the other four cannot: it \emph{marches}. Each
timestep begins exactly where the last one ended, so the solves form a chain, a
\texttt{fold}, that no machine on earth can run out of order. This is the one
genuinely sequential driver of the simulator. We assemble the three operators
$\Kop,\Cop,\Mop$ once, build an implicit ODE integrator once, and then thread a
single field-state forward through time, one opaque step at a time---and out
falls not a matrix but a \emph{movie}, the whole time history of the field. The
machinery of one step we built in \cref{ch:time}; here we compose it into the
driver and draw, as vividly as we can, the contrast between the map drivers and
this fold.}

\section{The physical question}
\label{sec:tr-question}

An engineer hands us a device and a \emph{pulse}: ``inject this current burst and
show me what happens.'' Not a steady voltage, not a single tone---a short,
broadband packet of current $\vJ(t)$ that arrives, excites the structure, and
rings down. The question is not ``what is the capacitance'' or ``what is the
response at \SI{5}{GHz}''; it is \emph{``what does the field do, frame by frame,
as time runs?''} The answer is a time history: the fields, the port voltages and
currents, the stored energy, each sampled at every instant of a simulated
interval. A movie, not a snapshot (\cref{fig:tr-setup}).

\begin{physicsbox}
A current pulse $\vJ(t)$---a modulated wave packet, a step, an impulse---is
injected into a device at $t=t_0$. The fields it launches propagate, reflect off
boundaries, couple into ports, and decay as energy bleeds away through loss and
radiation. At each instant $t_n$ the simulation holds a complete field-state: the
electric field $\vE(\vx,t_n)$ and magnetic field $\vB(\vx,t_n)$ everywhere in the
domain. The output is the whole sequence
$\vE(\cdot,t_0),\vE(\cdot,t_1),\vE(\cdot,t_2),\dots$---a time-march filmstrip, one
frame per timestep. Nothing here is frozen ($\partial_t\neq0$) and nothing is
assumed single-frequency: time is kept in full.
\end{physicsbox}

\begin{figure}[t]
\centering
\begin{tikzpicture}[font=\footnotesize, scale=1.0]
  % the drive pulse on the left
  \begin{scope}[shift={(-4.4,0)}]
    \draw[->, simGray] (-0.3,0)--(1.7,0) node[right,font=\scriptsize]{$t$};
    \draw[->, simGray] (0,-0.5)--(0,0.95);
    \draw[simRust, very thick, domain=-0.2:1.5, samples=60, smooth]
      plot (\x, {0.75*exp(-9*(\x-0.55)^2)*cos(deg(13*(\x-0.55)))});
    \node[font=\scriptsize, text=simRust] at (0.7,-0.85) {pulse $\vJ(t)$};
  \end{scope}
  \draw[flow] (-2.5,0) -- (-1.7,0);
  % three field snapshots t0, t1, t2 ...
  \foreach \k/\dx/\amp in {0/0/0.15, 1/2.4/0.55, 2/4.8/0.32}{
    \begin{scope}[shift={(\dx,0)}]
      \draw[simGray, semithick, rounded corners=2pt] (-0.85,-1.05) rectangle (0.85,1.05);
      % wavy field inside, amplitude grows then shrinks
      \draw[simBlue, thick, domain=-0.8:0.8, samples=40, smooth]
        plot (\x, {\amp*sin(deg(7*\x))});
      \draw[simBlue, thick, domain=-0.8:0.8, samples=40, smooth]
        plot (\x, {-\amp*sin(deg(7*\x))});
      \node[font=\scriptsize, anchor=north] at (0,-1.1) {$t_\k$};
    \end{scope}}
  \node[font=\scriptsize] at (6.1,0) {$\cdots$};
  \node[font=\scriptsize\itshape, text=simGreen] at (2.4,1.5)
    {the field evolves: a frame per timestep};
\end{tikzpicture}
\caption[The transient setup---a pulse and a filmstrip]{The transient analysis
in one picture. A current pulse $\vJ(t)$ (rust, left) is injected into the device;
the fields it launches evolve in time, and the simulation records a complete
field-state snapshot $(\vE,\vB)$ at every timestep $t_0,t_1,t_2,\dots$ (blue, the
field amplitude rising as the pulse arrives and decaying as energy leaves). The
product is not a single number but the \emph{whole filmstrip}---the time history
of the field. This is the analysis of choice for broadband excitation, for
nonlinear materials, or simply for watching a wave propagate.}
\label{fig:tr-setup}
\end{figure}

When do you reach for it? Three situations recommend the transient analysis over
its frequency-domain cousins. First, \emph{broadband excitation}: a single pulse
contains a whole band of frequencies at once (a narrow pulse in time is a wide
band in frequency), so one transient run can recover the response across an entire
band---where the driven analysis of \cref{ch:driven} sweeps one frequency per
solve. Second, \emph{nonlinearity}: a material whose properties depend on the
field amplitude has no single-frequency steady state to solve for, and only a
time-march can follow it. Third, simply \emph{watching a wave move}: when the
physical insight you want is spatial and temporal---a pulse crossing a junction, a
mode building up in a cavity---the movie \emph{is} the answer.

The transient analysis is the time-domain complement of the driven analysis. The
two compute the same physics by two routes of the harmonic rule
(\cref{ch:reductions-pde}): the driven analysis works one frequency at a time; the
transient analysis works all frequencies at once, through a pulse, and recovers
the spectrum afterward by a Fourier transform. We will close the chapter on
exactly that bridge (\S\ref{sec:tr-fourier}).

\section{Assemble, once: the three operators and the integrator}
\label{sec:tr-assemble}

The transient pipeline runs the same skeleton as the anchor
(\cref{ch:situating}),
\[
  \texttt{config} \to \underbrace{\texttt{assemble}}_{\S\ref{sec:tr-assemble}}
  \to \underbrace{\texttt{solve}}_{\S\ref{sec:tr-solve}}
  \to \underbrace{\texttt{reduce}}_{\S\ref{sec:tr-reduce}} \to \texttt{product},
\]
and the assemble stage, as always, runs \emph{once}, before any marching begins.
What differs from electrostatics is only \emph{how many} operators it builds and
\emph{what} it hands them to.

\subsection{Three operators on \texorpdfstring{$\Hcurl$}{H(curl)}}

The governing physics is the time-domain wave equation we derived in
\cref{ch:reductions-pde} (\S\ref{sec:transient}). Discretizing it in space only
(the method of lines) turns Maxwell's time-domain system into a \emph{second-order}
semi-discrete system of ordinary differential equations in the
degree-of-freedom vector $\vs(t)$:
\begin{equation}
  \boxed{\;\Mop\,\ddot\vs + \Cop\,\dot\vs + \Kop\,\vs = \vJ(t)\;}
  \label{eq:tr-secondorder}
\end{equation}
the damped-oscillator system of \cref{eq:transient-ode}. Here $\Mop$ is the
$\varepsilon$-weighted \emph{mass}, $\Cop$ the $\sigma$-weighted \emph{damping},
$\Kop$ the curl--curl \emph{stiffness}, and $\vJ(t)$ the discretized forcing
$-\partial_t\vJ_{\text{src}}$. The unknown is the electric field, a vector field,
so it lives in the Nédélec edge space $\Hcurl$ of \cref{ch:functionspaces}---the
same space as the driven, magnetostatic, and eigenmode analyses, and the swap from
electrostatics' scalar $\Hone$ that \cref{tab:es-swaps} recorded.

Where electrostatics assembled one operator, transient assembles \emph{three}, by
three calls to the same assembly fold \code{fe\_assemble} (\cref{ch:fem}), each
over its own one-term weak form:
\begin{lstlisting}[style=l4style]
-- assemble the three time-domain operators, each a one-term fe_assemble fold; ALL once
space = nd_space cfg                              -- the Nedelec H(curl) FE space (read-only)
k = fe_assemble space [ curl_curl  (reluctivity   cfg) ]  -- (1) stiffness K  (nu * curl)
c = fe_assemble space [ damping     (conductivity  cfg) ]  --     damping  C  (sigma * .)
m = fe_assemble space [ mass        (permittivity  cfg) ]  --     mass     M  (epsilon * .)
\end{lstlisting}
This is the $\Kop/\Mop/\Cop$ recurrence of \cref{ch:reductions-pde} made concrete:
the three operators are the fixed, reusable raw material, and every analysis is a
recipe for recombining them. The transient recipe is
\cref{eq:tr-secondorder}.\footnote{In Palace the three matrices are assembled once
inside the \code{TimeOperator} constructor:
\code{K = space\_op.GetStiffnessMatrix(...)},
\code{C = space\_op.GetDampingMatrix(...)},
\code{M = space\_op.GetMassMatrix(...)}
(\code{palace/models/timeoperator.cpp:65-67}). All three are built before the time
loop and held unchanged for the entire march.}

\subsection{Build the integrator once, and seed the state}

The assemble stage does two more things, both once. First it constructs the
\emph{ODE integrator}---the routine that will advance the state one step at a time.
As \cref{ch:time} explained at length, the semi-discrete system
\cref{eq:tr-secondorder} is stiff (its largest eigenvalue is set by the smallest
mesh element), so an explicit integrator would be pinned to absurdly tiny steps,
and Palace integrates \emph{implicitly}. It recasts \cref{eq:tr-secondorder} as a
first-order system in the doubled state $\vy=(\vs,\dot\vs)$ and hands it to a
library integrator---generalized-$\alpha$, SDIRK, or an ARKStep---chosen by a
config token and constructed \emph{once}, outside the loop
(\cref{sec:integrators}).\footnote{Palace builds a
\code{TimeDependentFirstOrderOperator} from $\Kop,\Cop,\Mop$ and constructs the
integrator at \code{palace/models/timeoperator.cpp:311-335}, before the driver's
time loop at \code{palace/drivers/transientsolver.cpp:77}.} Second it seeds the
initial field-state $\vy_0$---typically rest, $\vs_0=\dot\vs_0=\mathbf0$---the value
the fold will thread forward.

\begin{keyidea}
The transient assemble stage builds, \emph{once}, everything the march will reuse:
the three operators $\Kop,\Cop,\Mop$ on $\Hcurl$ (by three \code{fe\_assemble}
folds), the implicit ODE integrator (constructed from those operators, the step
machinery of \cref{ch:time}), and the seed field-state $\vy_0$. None of these
changes during the march---the operator capture is hoisted out of the time loop,
exactly as the electrostatic operator was hoisted out of the terminal loop
(\cref{ch:electrostatics}). The hoist is the same; what differs is the loop it sits
outside of.
\end{keyidea}

\section{Solve: the state-threaded fold}
\label{sec:tr-solve}

Here is the corner this chapter exists to teach. With the operators assembled, the
integrator built, and the state seeded, the solve stage marches. And the march is
\emph{not} a map. It is a \texttt{fold}: the one solve corner of \cref{ch:situating}
in which each solve begins from the previous solve's output. This is the structural
centerpiece of the chapter, and the rest of Part~V has been pointing at it.

\subsection{The crucial contrast: a fold, not a map}

Recall the fundamental split of \cref{ch:combinators}, which the anchor chapter
leaned on (\cref{fig:es-solvefamily}). A \texttt{map} runs a family of
\emph{independent} computations: no carry passes between elements, so they
reorder freely and run in parallel. A \texttt{fold} threads a \emph{carry} through
a chain: each step consumes what the last produced, so the steps are strictly
ordered. The electrostatic and driven solve corners are maps---a fan of
independent solves over a family of terminals or frequencies. The transient solve
corner is a fold, and the difference is the whole point.

In the transient march the field-state at time $t_{n+1}$ is computed \emph{from}
the field-state at time $t_n$ by one timestep. Step $n+1$ cannot begin until step
$n$ has finished, because step $n$'s output \emph{is} step $n+1$'s input. The state
is threaded forward along a single timeline. You cannot solve the timesteps in
parallel, in any order, or on separate machines, because there is no ``step 5'' to
hand a machine until ``step 4'' has produced the state it starts from. This is the
\emph{state-threaded fold}---the combinator \code{fold\_solve}---and the transient
march is its purest and primary witness. \Cref{fig:tr-foldvmap} sets it against the
electrostatic map, the two side by side.

\begin{figure}[p]
\centering
\begin{tikzpicture}[font=\footnotesize]
  % ===== LEFT: the map drivers (electrostatic / driven) =====
  \begin{scope}
    \node[font=\small\bfseries, text=simBlue] at (2.1,3.0)
      {\textsc{map} drivers: independent, parallel};
    \node[font=\scriptsize\itshape, text=simBlue] at (2.1,2.6)
      {electrostatic, driven};
    \node[opbox, minimum width=13mm, fill=simBgBlue, draw=simBlue] (K) at (2.1,1.7)
      {operator};
    \foreach \i/\x in {1/0.3,2/1.5,3/2.7,4/3.9}{
      \node[data, minimum width=7mm, minimum height=5mm] (b\i) at (\x,0.75) {$\vb_{\i}$};
      \node[product, minimum width=7mm, minimum height=5mm] (x\i) at (\x,-0.3) {$\vx_{\i}$};
      \draw[flow, shorten >=1pt] (K.south) -- (b\i.north);
      \draw[flow, shorten >=1pt] (b\i.south) -- (x\i.north);
    }
    \node[font=\scriptsize, text=simGreen, align=center] at (2.1,-1.1)
      {NO arrows between solves:\\ reorderable, embarrassingly parallel};
  \end{scope}
  \draw[simGray!55, very thick] (4.9,-1.5) -- (4.9,3.2);
  % ===== RIGHT: the transient fold =====
  \begin{scope}[xshift=5.6cm]
    \node[font=\small\bfseries, text=simRust] at (2.4,3.0)
      {\textsc{fold} driver: chained, sequential};
    \node[font=\scriptsize\itshape, text=simRust] at (2.4,2.6)
      {transient (the one fold driver)};
    \node[data, minimum width=8mm] (y0) at (-0.25,1.5) {$\vy_0$};
    \foreach \i/\x in {1/1.0,2/2.5,3/4.0}{
      \node[opbox, minimum width=11mm] (g\i) at (\x,1.5) {step};
    }
    \node[data, minimum width=8mm] (y3) at (5.25,1.5) {$\vy_3$};
    \draw[flow] (y0.east) -- (g1.west);
    \draw[flow] (g1.east) -- node[above,font=\scriptsize]{$\vy_1$} (g2.west);
    \draw[flow] (g2.east) -- node[above,font=\scriptsize]{$\vy_2$} (g3.west);
    \draw[flow] (g3.east) -- (y3.west);
    % each step hides a ksp_solve
    \foreach \i/\x in {1/1.0,2/2.5,3/4.0}{
      \node[extern, minimum width=11mm, minimum height=5mm, font=\scriptsize]
        (k\i) at (\x,0.3) {\code{ksp\_solve}};
      \draw[refedge] (g\i.south) -- (k\i.north);
    }
    \node[font=\scriptsize, text=simRust, align=center] at (2.4,-0.85)
      {horizontal carry $\vy_n\!\to\!\vy_{n+1}$:\\ strictly ordered---cannot overlap timesteps};
  \end{scope}
\end{tikzpicture}
\caption[The map drivers versus the transient fold]{The centerpiece contrast. On
the \emph{left}, the map drivers (electrostatic, driven): one operator feeds a fan
of independent solves $\vb_i\mapsto\vx_i$, with \emph{no} arrows between them---they
are reorderable, parallelizable, embarrassingly parallel over the family. On the
\emph{right}, the transient fold: the doubled state $\vy_0\to\vy_1\to\vy_2\to\vy_3$
is threaded \emph{forward}, each step consuming the prior state, so the steps are
strictly ordered and \emph{cannot} overlap. The horizontal carry arrows---present on
the right, absent on the left---are the whole difference: \code{fold\_solve} versus
\code{solve\_family}. Each fold step \emph{additionally} hides an implicit
\code{ksp\_solve} (\cref{ch:krylov}), but the implicit solve is not what makes it a
fold; the \emph{carry} is.}
\label{fig:tr-foldvmap}
\end{figure}

\begin{keyidea}
\textbf{Transient is the one sequential driver: a fold where each step begins where
the last ended---the opposite of the map drivers.} Electrostatic and driven solve a
\emph{family} of independent problems (a \texttt{map}: reorderable, parallel over
terminals or frequencies). Transient threads a \emph{single field-state} forward
through time (a \texttt{fold}: each step's input is the prior step's output). The
presence or absence of the carry between solves is the entire structural
difference, and it is the difference between a contact sheet of independent
snapshots and a movie in which every frame depends on the last. No amount of
hardware can parallelize the transient march across timesteps, because step $n+1$
literally has nothing to compute until step $n$ is done.
\end{keyidea}

\subsection{\texttt{fold\_solve}, in the calculus}

In the vocabulary of \cref{ch:combinators}, the march is one expression. Capture
the operator once, seed the state once, and \code{foldl} the opaque per-step
operator over the schedule of times:
\begin{lstlisting}[style=l4style]
-- the state-threaded march: foldl the per-step operator over the fixed schedule
fold_solve :: OpParams -> TimeState -> [Time] -> TimeState
fold_solve op s0 schedule = foldl (\s t -> time_step_op op s t) s0 schedule

-- the OPAQUE per-step operator the fold quantifies over:
-- advances the field-state one step; bottoms out in the library integrator
time_step_op :: OpParams -> TimeState -> Time -> TimeState
\end{lstlisting}
Read the signature and the whole structure of the march is in it. The seed
\code{s0} and the result are both a \code{TimeState}---the field-state carried
forward. The operator \code{op} (the captured $\Kop,\Cop,\Mop$ and the chosen
integrator) is bound \emph{once}, before the fold, and threaded unchanged into every
step: the construction hoist of \S\ref{sec:tr-assemble} made structural by the type.
And the step \lstinline[style=l4style]|time_step_op op s t| \emph{reads} \code{s},
the prior carry---\emph{that read} is what makes this a fold and forbids reordering
at the type level. A map's step would ignore the accumulator; a fold's reads it.

\code{time\_step\_op} is deliberately \emph{opaque}. It bottoms out in the library
integrator's one step---an MFEM \code{ODESolver::Step} call---which internally does
the implicit linear solve that \cref{ch:time} dissected
(\cref{fig:tr-innersolve}).\footnote{The fold body is \code{time\_op.Step(t,
delta\_t)} (\code{palace/drivers/transientsolver.cpp:93}), which dispatches to
\code{ode->Step(sol, t, dt)} (\code{palace/models/timeoperator.cpp:410})---the
opaque MFEM step advancing the persistent \code{sol} field-state in place. The prior
step's \code{sol} is the next step's input: the genuine \code{foldl}.} The
combinator only quantifies over that step; it does not look inside it.
\code{fold\_solve} owns the loop, and the library owns one step.

\begin{figure}[t]
\centering
\begin{tikzpicture}[font=\footnotesize]
  % the fold spine
  \node[data, minimum width=8mm] (sn) at (0,0) {$\vy_n$};
  \node[opbox, minimum width=26mm, minimum height=16mm, fill=simBgRust, draw=simRust,
    align=center] (step) at (3.0,0) {\code{time\_step\_op}\\[1pt]\scriptsize one implicit step};
  \node[data, minimum width=10mm] (sn1) at (6.2,0) {$\vy_{n+1}$};
  \draw[flow] (sn) -- (step);
  \draw[flow] (step) -- (sn1);
  % the inner ksp_solve, blown up below
  \node[extern, minimum width=30mm, minimum height=10mm, align=center]
    (ksp) at (3.0,-2.2) {\scriptsize MFEM \code{ODESolver::Step}\\[-1pt]\scriptsize $(\mat I-\Delta t\,\mat J)\,\vy_{n+1}=\text{rhs}$\\[-1pt]\scriptsize one \code{ksp\_solve}};
  \draw[refedge] (step.south) -- (ksp.north)
    node[midway, right, font=\scriptsize]{hides};
  \node[font=\scriptsize\itshape, text=simGray, align=left, anchor=west] at (-1.6,-2.2)
    {\emph{inside} each};
  \node[font=\scriptsize\itshape, text=simGray, align=left, anchor=west] at (-1.6,-2.55)
    {opaque step:};
\end{tikzpicture}
\caption[The implicit solve inside each fold step]{Each fold step is opaque, but it
is not empty: inside \code{time\_step\_op} the integrator solves one linear system.
Because the integrator is \emph{implicit} (the cure for the stiffness of
\cref{ch:time}), advancing $\vy_n\to\vy_{n+1}$ requires solving
$(\mat I-\Delta t\,\mat J)\,\vy_{n+1}=\text{rhs}$ for the unknown next state---a
single \code{ksp\_solve}, the Krylov solve of \cref{ch:krylov}, nested in the march.
The fold quantifies over this step; the library performs it. Note the two distinct
costs the fold carries: the \emph{carry} (which makes the march sequential) and the
\emph{implicit solve} (which makes each step expensive). They are independent---a
fold could be cheap per step, and a map could be expensive per step---but the
transient march has both.}
\label{fig:tr-innersolve}
\end{figure}

\subsection{The pipeline in one line}

The two stages compose into a single expression. At the framework level the whole
transient feature is
\begin{lstlisting}[style=l4style]
transient = fold_solve . fe_assemble
\end{lstlisting}
read right to left: \code{fe\_assemble} builds the operators once, and
\code{fold\_solve} threads the time-march fold over them. Spelled out with the data
flowing through it,
\begin{lstlisting}[style=l4style]
-- inputs = config; output = the time-domain field-state trajectory (the product)
transient :: TransientConfig -> FieldTrajectory
transient cfg =
  let space    = nd_space cfg                             -- the H(curl) FE space
      (k,c,m)  = ( fe_assemble space [ curl_curl (reluctivity cfg) ]   -- (1) K, C, M
                 , fe_assemble space [ damping    (conductivity cfg) ]  --     assembled
                 , fe_assemble space [ mass       (permittivity cfg) ] )--     ONCE
      op       = time_operator (k,c,m) (dJdt cfg)         -- captured ODE operator (read-only)
      s0       = init_state cfg                           -- seed field-state (rest)
      schedule = uniform_steps (delta_t cfg) (n_step cfg) -- the FIXED [Time] schedule
  in  fold_solve op s0 schedule                           -- (2) the state-threaded march FOLD
\end{lstlisting}
\Cref{fig:tr-pipeline} draws it. Set it beside the anchor diagram
(\cref{fig:es-pipeline}) and the swap is exactly the one \cref{tab:es-swaps}
promised: the FE space moves from $\Hone$ to $\Hcurl$, the solve corner moves from
the \code{solve\_family} \emph{map} to the \code{fold\_solve} \emph{fold}, and the
reduction moves from the Gram to a trajectory table. The skeleton is identical; one
corner is swapped.

\begin{figure}[p]
\centering
\begin{tikzpicture}[node distance=6mm and 6mm, font=\footnotesize]
  \node[data, align=center, text width=18mm] (cfg)
    {config\\[1pt]\footnotesize $\varepsilon,\sigma,\nu$,\\mesh, pulse};
  \node[stage, right=of cfg, align=center, text width=21mm] (asm)
    {\textbf{assemble}\\[1pt]\footnotesize \code{fe\_assemble}\\$\Kop,\Cop,\Mop$ (once)};
  \node[stage, right=of asm, align=center, text width=22mm, fill=simBgRust, draw=simRust]
    (solve) {\textbf{solve}\\[1pt]\footnotesize \code{fold\_solve}\\(time march)};
  \node[stage, right=of solve, align=center, text width=18mm, fill=simBgGold, draw=simGold]
    (red) {\textbf{reduce}\\[1pt]\footnotesize per-step\\readout};
  \node[product, right=of red, align=center, text width=17mm] (prod)
    {trajectory\\[1pt]\footnotesize field history};
  \draw[flow] (cfg) -- (asm);
  \draw[flow] (asm) -- node[above,font=\scriptsize]{$\Kop,\Cop,\Mop$} (solve);
  \draw[flow] (solve) -- node[above,font=\scriptsize]{$\{\vy_n\}$} (red);
  \draw[flow] (red) -- (prod);
  % the self-loop on the solve box: the carry
  \draw[backflow, shorten >=1pt] (solve.north) ++(-0.4,0)
    .. controls +(0,7mm) and +(0,7mm) .. ($(solve.north)+(0.4,0)$)
    node[midway, above, font=\scriptsize, text=simRust]
      {the carry $\vy_n\!\to\!\vy_{n+1}$};
  % corner annotations
  \node[font=\scriptsize\itshape, text=simRust, below=9mm of solve, align=center]
    {state-threaded fold\\(\cref{ch:situating}, corner 3)};
  \node[font=\scriptsize\itshape, text=simGold, above=8mm of red, align=center]
    {trajectory, not a matrix};
\end{tikzpicture}
\caption[The transient pipeline]{The complete transient pipeline, the composition
$\code{fold\_solve}\circ\code{fe\_assemble}$. A configuration of
materials, mesh, and excitation pulse enters; \code{fe\_assemble} builds the three
operators $\Kop,\Cop,\Mop$ \emph{once} (and, with them, the integrator and seed
state); \code{fold\_solve} threads the time-march fold, producing the field-state
trajectory $\{\vy_n\}$; and a per-step readout reduces each state to the physical
products. The rust self-loop on the solve box is the feature special to this
analysis: the \emph{carry}, $\vy_n\to\vy_{n+1}$, that makes the march sequential.
Compared with the anchor (\cref{fig:es-pipeline}), exactly one corner is swapped:
the \code{solve\_family} map becomes the \code{fold\_solve} fold.}
\label{fig:tr-pipeline}
\end{figure}

\subsection{Checkpoint and resume: the schedule-split property}

The fold's sequentiality sounds like pure cost, but it comes with a genuine
compensation that the map cannot offer: \emph{restartability}. A left fold over a
concatenated schedule satisfies a clean algebraic law (\cref{ch:time},
\cref{eq:schedule-split}). Split the schedule into a prefix $a$ and a suffix $b$;
then folding the whole is the same as folding the prefix and folding the suffix
\emph{from the state the prefix produced}:
\begin{equation}
  \code{fold\_solve}\ \code{op}\ \vy_0\ (a \mathbin{+\!\!+} b)
  \;=\;
  \code{fold\_solve}\ \code{op}\ \bigl(\code{fold\_solve}\ \code{op}\ \vy_0\ a\bigr)\ b .
  \label{eq:tr-split}
\end{equation}
Read operationally, this says: run the prefix, save the final state $\vy_{|a|}$ to
disk, stop; later, load it and run the suffix seeded from it. The law guarantees the
result is the \emph{identical} trajectory---bit for bit---as the uninterrupted run,
because the only thing that crosses the split point is the carry, and the carry is
exactly what we saved (\cref{fig:tr-checkpoint}). Long runs survive faults; expensive
runs split across job-scheduler windows; none of it perturbs the answer. The same
single threaded carry that \emph{forbids} parallelism across timesteps is what makes
a checkpoint a single saved value rather than a replay of the whole history. A map
parallelizes but has no notion of ``resume''---there is no carry to save.

\begin{figure}[t]
\centering
\begin{tikzpicture}[font=\footnotesize]
  % top: uninterrupted
  \node[font=\scriptsize\bfseries, text=simInk, anchor=west] at (-0.3,1.9)
    {one uninterrupted march};
  \node[data, minimum width=8mm] (u0) at (0,1.2) {$\vy_0$};
  \foreach \i/\x in {1/0.95,2/1.95,3/2.95,4/3.95}{
    \node[opbox, minimum width=7mm, minimum height=5mm] (us\i) at (\x,1.2) {};}
  \node[data, minimum width=8mm] (u4) at (4.95,1.2) {$\vy_4$};
  \draw[flow] (u0)--(us1);\draw[flow] (us1)--(us2);\draw[flow] (us2)--(us3);
  \draw[flow] (us3)--(us4);\draw[flow] (us4)--(u4);
  \node[font=\scriptsize, text=simGray] at (2.5,0.72) {schedule $a \mathbin{+\!\!+} b$};
  % bottom: split
  \node[font=\scriptsize\bfseries, text=simRust, anchor=west] at (-0.3,-0.1)
    {split: run $a$, save carry, resume $b$};
  \node[data, minimum width=8mm] (b0) at (0,-0.8) {$\vy_0$};
  \foreach \i/\x in {1/0.95,2/1.95}{
    \node[opbox, minimum width=7mm, minimum height=5mm] (bs\i) at (\x,-0.8) {};}
  \node[product, minimum width=9mm, fill=simBgGold, draw=simGold] (ckpt) at (2.95,-0.8) {$\vy_2$};
  \foreach \i/\x in {3/3.95,4/4.95}{
    \node[opbox, minimum width=7mm, minimum height=5mm] (bs\i) at (\x,-0.8) {};}
  \node[data, minimum width=8mm] (b4) at (5.95,-0.8) {$\vy_4$};
  \draw[flow] (b0)--(bs1);\draw[flow] (bs1)--(bs2);\draw[flow] (bs2)--(ckpt);
  \draw[flow] (ckpt)--(bs3);\draw[flow] (bs3)--(bs4);\draw[flow] (bs4)--(b4);
  \node[font=\scriptsize, text=simRust] at (1.4,-1.3) {fold $a$};
  \node[font=\scriptsize, text=simRust] at (4.4,-1.3) {fold $b$ from $\vy_2$};
  \node[font=\scriptsize, text=simGold!75!black, align=center] at (2.95,-1.75)
    {save / load (checkpoint)};
  \node[font=\Large, text=simGray] at (5.6,0.2) {$=$};
\end{tikzpicture}
\caption[Checkpoint and resume as a schedule split]{The schedule-split law
\cref{eq:tr-split} drawn as checkpoint-and-resume. \emph{Top}: the uninterrupted
march folds the whole schedule $a\mathbin{+\!\!+}b$ from $\vy_0$ to $\vy_4$.
\emph{Bottom}: the same march split at the midpoint---fold the prefix $a$ to the
carry $\vy_2$, \emph{save it} (the gold checkpoint), stop, then later load $\vy_2$
and fold the suffix $b$ from it. The two produce the identical $\vy_4$, because the
only thing crossing the cut is the carry. Restartability is not bolted onto the
fold; it \emph{is} what the fold structure means---a clean algebraic property, not
an engineering approximation.}
\label{fig:tr-checkpoint}
\end{figure}

\section{Reduce: reading the trajectory}
\label{sec:tr-reduce}

The solve stage has handed us not a family of fields but a \emph{trajectory}: the
field-state at every step of the march. The reduce stage reads the physical products
from it. Where the electrostatic reduction collapsed a family into a single
capacitance matrix (\cref{ch:electrostatics}), the transient ``reduction'' is
mostly a per-step \emph{readout}---it runs once per frame, and the output is the
whole time history, not one number.

\subsection{What the driver reads at each step}

At each step the transient driver recovers three kinds of product from the current
field-state:

\begin{description}[style=nextline, leftmargin=0pt, font=\normalfont\itshape]
  \item[The fields.] From each step's state the electric and magnetic fields
  $(\vE,\vB)$ are recovered (the \code{GetE}/\code{GetB} of the time operator) and
  written out---the ``movie'' frames, one per timestep.\footnote{The per-step readout
  is \code{time\_op.GetE()} / \code{time\_op.GetB()}
  (\code{palace/drivers/transientsolver.cpp:98-99}), consumed immediately by the
  per-step postprocess \code{post\_op.MeasureAndPrintAll(...)}
  (\code{:104}).}
  \item[Port voltages and currents.] Integrating the fields against the port modes
  gives a time-domain voltage and current at each terminal: the waveforms an
  engineer reads off an oscilloscope, computed step by step.
  \item[Energy versus time.] The field energy, evaluated at each step, traces how
  energy enters with the pulse, sloshes between electric and magnetic forms, and
  bleeds away through the damping $\Cop$ and the ports. This is the
  \emph{driver-agnostic} energy reduction of \cref{ch:reductions}---the same
  field-energy form the electrostatic diagonal used (\cref{ch:electrostatics}),
  here evaluated once per frame instead of once per terminal.
\end{description}

\Cref{fig:tr-energy} draws a representative trajectory: an injected pulse rings up
the stored energy, which then decays as loss and radiation carry it off. Because
each curve is sampled once per fold step, the resolution of the readout is exactly
the resolution of the march---one frame per timestep.

\begin{figure}[t]
\centering
\begin{tikzpicture}
\begin{axis}[
    width=0.80\linewidth, height=50mm,
    xlabel={time $t$}, ylabel={},
    xmin=0, xmax=10, ymin=-1.15, ymax=1.30,
    axis lines=left, xtick={0,2,4,6,8,10}, ytick={-1,0,1},
    yticklabel style={font=\scriptsize}, xticklabel style={font=\scriptsize},
    xlabel style={font=\scriptsize},
    legend style={font=\scriptsize, draw=simGray!50, at={(0.98,0.98)},
      anchor=north east, legend cell align=left},
    every axis plot/.append style={very thick, samples=120, domain=0:10},
  ]
  \addplot[simRust] {0.9*exp(-2*(x-1.5)^2)*cos(deg(6*(x-1.5)))};
  \addlegendentry{drive $\vJ(t)$ (a pulse)}
  \addplot[simBlue] {0.8*(1-exp(-3*x))*exp(-0.35*x)*cos(deg(4*x-1))};
  \addlegendentry{field $E(t)$ (rings, decays)}
  \addplot[simGreen, densely dashed] {0.95*(1-exp(-2.2*x))*exp(-0.5*x)};
  \addlegendentry{stored energy (rise then decay)}
\end{axis}
\end{tikzpicture}
\caption[A transient field and energy trajectory]{A representative readout of the
transient march. The pulse drive $\vJ(t)$ (rust) injects a short burst; a field
component $E(t)$ (blue) rings up as the pulse arrives and decays as energy leaves;
the stored energy (green, dashed) rises with the injection and decays monotonically
thereafter, bled away by the damping $\Cop$ and carried off through the ports. Each
curve is sampled once per fold step: the readout resolution is the march resolution,
one frame per timestep. The output is this \emph{whole history}, not a single
reduced number.}
\label{fig:tr-energy}
\end{figure}

\section{The bridge to the frequency domain}
\label{sec:tr-fourier}

One readout deserves its own section, because it ties the transient analysis back
to the driven analysis of \cref{ch:driven} and closes a loop the whole of Part~V has
left open. A transient run injects a \emph{broadband} pulse---a superposition of
many frequencies at once. Record the input waveform and the corresponding output
waveform in time, take the Fourier transform of each, and divide: the quotient is
the frequency response across the \emph{whole band} of the pulse, recovered \emph{in
a single run}.

This is the time-domain route to the same S-parameters the driven analysis computes
frequency by frequency. Where the driven analysis sweeps a \code{frequency\_sweep}
map---one frequency per solve, an independent family (\cref{ch:driven})---the
transient analysis folds one march and Fourier-transforms the result. The two compute
the same physics by the two routes of the harmonic rule
(\cref{ch:reductions-pde}): one frequency at a time, or all frequencies at once
through a pulse. \Cref{fig:tr-fourier} shows a sampled port-voltage signal beside its
transform, a resonance appearing as a peak in the spectrum.

\begin{figure}[t]
\centering
\begin{tikzpicture}
% LEFT: time-domain port signal
\begin{axis}[
    name=tplot,
    width=0.46\linewidth, height=44mm,
    xlabel={time $t$}, ylabel={port voltage $v(t)$},
    xmin=0, xmax=12, ymin=-1.1, ymax=1.1,
    axis lines=left, xtick={0,4,8,12}, ytick={-1,0,1},
    yticklabel style={font=\scriptsize}, xticklabel style={font=\scriptsize},
    xlabel style={font=\scriptsize}, ylabel style={font=\scriptsize},
    title style={font=\scriptsize\itshape}, title={time domain (the march)},
    every axis plot/.append style={very thick, samples=160, domain=0:12},
  ]
  % a decaying ring at a single frequency: the "resonance"
  \addplot[simBlue] {exp(-0.32*x)*cos(deg(5*x))};
\end{axis}
% RIGHT: frequency-domain magnitude with a peak
\begin{axis}[
    at={(tplot.right of east)}, anchor=left of west, xshift=10mm,
    width=0.46\linewidth, height=44mm,
    xlabel={frequency $f$}, ylabel={$|\hat v(f)|$},
    xmin=0, xmax=2, ymin=0, ymax=1.15,
    axis lines=left, xtick={0,0.5,1,1.5,2}, ytick={0,0.5,1},
    yticklabel style={font=\scriptsize}, xticklabel style={font=\scriptsize},
    xlabel style={font=\scriptsize}, ylabel style={font=\scriptsize},
    title style={font=\scriptsize\itshape}, title={frequency domain (FFT)},
    every axis plot/.append style={very thick, samples=160, domain=0:2},
  ]
  % a Lorentzian peak at f0 ~ 0.8 (= 5/(2 pi)), the spectrum of the decaying ring
  \addplot[simRust] {1/sqrt(1 + 80*(x-0.8)^2)};
  \draw[simGray, densely dashed] (axis cs:0.8,0) -- (axis cs:0.8,1.0);
  \node[font=\scriptsize, text=simGray, anchor=south] at (axis cs:0.8,1.0) {$f_0$};
\end{axis}
\end{tikzpicture}
\caption[A port signal and its Fourier transform]{The time$\to$frequency bridge.
\emph{Left}: a sampled port-voltage waveform $v(t)$ from one transient march---a
decaying ring at the structure's resonant frequency. \emph{Right}: its Fourier
transform $|\hat v(f)|$, a peak at the resonance $f_0$ (a slow decay in time is a
sharp peak in frequency---a high-$Q$ resonance). Dividing the transform of the
\emph{output} by that of the \emph{input} pulse recovers the frequency response
across the whole band, in one run. This is the same frequency-domain information the
driven analysis (\cref{ch:driven}) computes one frequency at a time; the transient
fold plus an FFT is its broadband, all-at-once complement.}
\label{fig:tr-fourier}
\end{figure}

\section{Cost and stability}
\label{sec:tr-cost}

Two practical facts about the march, both carried over from the machinery of
\cref{ch:time}, are worth collecting before the worked examples.

\paragraph{Implicit stepping, and the timestep trade.} Each fold step hides an
implicit linear solve precisely because the semi-discrete system is \emph{stiff}: an
explicit step would be pinned to $\Delta t<2/\lambda_{\max}$ by the fastest mesh
mode, a mode physically irrelevant to the answer (\cref{ch:time}). An implicit step
is unconditionally stable, so $\Delta t$ is chosen for \emph{accuracy on the modes we
care about}, not forced small by modes we do not. The trade is direct: a smaller
$\Delta t$ resolves faster physics and reduces truncation error but costs more steps
(more \code{ksp\_solve}s); a larger $\Delta t$ is cheaper but coarser. The integrator's
order (\cref{ch:time}) sets how fast that error falls as $\Delta t$ shrinks.

\paragraph{Reuse the factorization across steps.} For a \emph{uniform} schedule the
step operator $\mat I-\Delta t\,\mat J$ is built from the captured $\Kop,\Cop,\Mop$
and the fixed $\Delta t$, none of which change from step to step---so it is a
\emph{constructed operator}, fixed for the whole march (\cref{ch:operators}). Its
factorization (or its preconditioner) is therefore computed \emph{once} and reused
every step, exactly as the electrostatic operator's factorization was reused across
the terminal family (\cref{ch:electrostatics}). The operator-capture-once hoist of
\code{fold\_solve} is what licenses it: \code{op} is bound before the fold, and every
step reads the same operator. This is the fold's version of the dividend the
fixed-operator map paid---``assemble once'' quietly licenses ``factor once'' here too,
even though the solves are chained rather than independent.

\begin{notationbox}
The transient march carries \emph{two} costs that are easy to conflate but logically
independent. The \emph{carry} makes the march \emph{sequential}: timesteps cannot
overlap, regardless of how cheap each one is. The \emph{implicit solve} makes each
step \emph{expensive}: one \code{ksp\_solve} per step, regardless of whether the
steps are chained. A hypothetical explicit transient march would still be a fold
(sequential) but with cheap steps; an electrostatic map is parallel but with
expensive steps. Transient is the corner that pays both: sequential \emph{and}
per-step-expensive. Keeping the two costs separate is the key to reasoning about it.
\end{notationbox}

\section{Two worked examples}
\label{sec:tr-worked}

The pipeline is general; we now run it on problems small enough to do by hand. The
first marches a two-degree-of-freedom oscillator a few implicit steps, showing each
step solve a small linear system and the state thread forward---and contrasts it
explicitly with the independent map a \code{solve\_family} would run. The second
takes a sampled port signal and shows its DFT recovering a resonance, making the
time$\to$frequency bridge of \S\ref{sec:tr-fourier} numerical.

\begin{workedexample}{Marching a two-DoF oscillator---a fold, not a map}{tr-twodof}
We run the real machinery---the recast, the per-step implicit solve, and the state
threading---on the smallest non-trivial system: a two-degree-of-freedom
mass--spring--damper $\Mop\ddot\vs+\Cop\dot\vs+\Kop\vs=\vJ(t)$, a hand-sized
stand-in for the semi-discrete transient system, with
\[
  \Mop=\begin{psmallmatrix}1&0\\0&1\end{psmallmatrix},\quad
  \Cop=\begin{psmallmatrix}0.2&0\\0&0.2\end{psmallmatrix},\quad
  \Kop=\begin{psmallmatrix}2&-1\\-1&2\end{psmallmatrix},\quad
  \vJ=\begin{psmallmatrix}1\\0\end{psmallmatrix}\ (\text{constant}),
\]
starting from rest, $\vs_0=\vv_0=(0,0)$. The off-diagonal of $\Kop$ couples the two
degrees of freedom.

\medskip
\emph{Recast and the step matrix.} With $\vv=\dot\vs$ the doubled state is
$\vy=(\vs,\vv)$. Take backward Euler with $\Delta t=0.1$. Substituting
$\vs_{n+1}=\vs_n+\Delta t\,\vv_{n+1}$ into the velocity update eliminates
$\vs_{n+1}$ and leaves a linear system for $\vv_{n+1}$ alone:
\[
  \underbrace{\bigl(\mat I+\Delta t\,\Cop+\Delta t^2\,\Kop\bigr)}_{\displaystyle\mat A}
  \,\vv_{n+1}=\vv_n+\Delta t\,(\vJ-\Kop\,\vs_n).
\]
This is the implicit step's linear solve---the $2\times2$ stand-in for the
\code{ksp\_solve} (\cref{fig:tr-innersolve}). With our numbers,
\[
  \mat A=\mat I+0.1\,\Cop+0.01\,\Kop
        =\begin{psmallmatrix}1.04&-0.01\\-0.01&1.04\end{psmallmatrix}.
\]
The matrix $\mat A$ is \emph{fixed for the whole march}---built from the captured
operators and $\Delta t$, none of which change---so it is factored \emph{once} and
reused every step (\S\ref{sec:tr-cost}). This is the operator-capture-once hoist of
\code{fold\_solve} in miniature: \code{op} is bound before the fold, and every step
reads the same $\mat A$.

\medskip
\emph{Thread the state forward.} Starting from $\vy_0=\mathbf0$, solve for $\vv_1$,
set $\vs_1=\Delta t\,\vv_1$, feed $\vy_1=(\vs_1,\vv_1)$ into the next step, and
repeat. Three steps:
\begin{center}
\setlength{\tabcolsep}{6pt}
\renewcommand{\arraystretch}{1.2}
\begin{tabular}{@{}clll@{}}
\toprule
step $n$ & RHS $\vv_n+\Delta t(\vJ-\Kop\vs_n)$ & solve $\vv_{n+1}=\mat A^{-1}(\cdot)$ & $\vs_{n+1}=\vs_n+\Delta t\,\vv_{n+1}$ \\
\midrule
0 & $(0.100,\ 0.000)$ & $(0.0961,\ 0.0009)$ & $(0.00961,\ 0.00009)$ \\
1 & $(0.1854,\ 0.0105)$ & $(0.1784,\ 0.0118)$ & $(0.02745,\ 0.00127)$ \\
2 & $(0.2589,\ 0.0264)$ & $(0.2492,\ 0.0278)$ & $(0.05237,\ 0.00405)$ \\
\bottomrule
\end{tabular}
\end{center}
Each row's input is the previous row's output---the carry $\vy_n\to\vy_{n+1}$
threaded forward, exactly the fold. The driven coordinate $s^{(1)}$ climbs toward
its static deflection while the coupled coordinate $s^{(2)}$ follows through the
off-diagonal stiffness; with the damping present, the march settles rather than
oscillating forever.

\medskip
\emph{Contrast: what a \code{solve\_family} map would do instead.} Suppose we had,
instead, three \emph{independent} right-hand sides $\vb_1,\vb_2,\vb_3$ and the one
fixed operator $\mat A$---an electrostatic-style \code{solve\_family}. We would solve
\[
  \mat A\,\vx_1=\vb_1,\qquad \mat A\,\vx_2=\vb_2,\qquad \mat A\,\vx_3=\vb_3,
\]
and \emph{nothing in $\vb_2$ depends on $\vx_1$}. We could compute $\vx_2$ first, or
on a separate machine, or all three at once: the solves commute. That is the map. The
transient table above is its exact negation: the RHS of step~1 is built from
$\vs_1$, which came from step~0; the RHS of step~2 from $\vs_2$, which came from
step~1. The carry column \emph{is} the dependence chain. Same fixed operator $\mat A$
in both; the difference is whether the right-hand sides are a fixed independent
\emph{family} (map) or a sequence each \emph{generated by the previous solve} (fold).
Three points are the chapter in miniature: \emph{(i)} the recast doubled
$\mathbb R^2$ into $\mathbb R^4$; \emph{(ii)} each step \emph{solved a linear system}
with $\mat A$ fixed and factored once; \emph{(iii)} the state was \emph{threaded}
from each step into the next---a \code{fold\_solve}, not a \code{solve\_family}. Were
we to checkpoint after step~1, saving $\vy_1$, and resume, the schedule-split law
\cref{eq:tr-split} guarantees the identical $\vy_3$.
\end{workedexample}

\begin{workedexample}{From a port waveform to a resonance: the DFT bridge}{tr-dft}
We make the time$\to$frequency bridge of \S\ref{sec:tr-fourier} numerical. Suppose a
transient march produced a sampled port-voltage waveform---a decaying ring at the
structure's resonant frequency. To keep the arithmetic in hand we sample $N=8$
points of a clean cosine ring,
\[
  v_n=\cos\!\Bigl(\tfrac{2\pi}{8}\,2\,n\Bigr),\qquad n=0,1,\dots,7,
\]
that is, a cosine that completes exactly \emph{two} full cycles over the eight
samples: $v=(1,0,-1,0,1,0,-1,0)$.

\medskip
\emph{The discrete Fourier transform.} The DFT of the sampled signal is
\[
  \hat v_k=\sum_{n=0}^{7} v_n\,e^{-i\,2\pi kn/8},\qquad k=0,1,\dots,7,
\]
each $\hat v_k$ measuring how much of the signal oscillates at frequency-index $k$
(i.e.\ at $k$ cycles over the window). Our signal is, by construction, a pure cosine
of \emph{two} cycles, so we expect all its weight to land on $k=2$ (and its mirror
$k=6=N-2$, since a real cosine splits its energy between $+f$ and $-f$). Compute the
suspect bins. For $k=2$, the exponential is $e^{-i\pi n/2}=(1,-i,-1,i,\dots)$, whose
\emph{real} part is $(1,0,-1,0,1,0,-1,0)$---identical to $v_n$---so
\[
  \hat v_2=\sum_{n}v_n\,e^{-i\pi n/2}
         =\underbrace{(1+1+1+1)}_{\text{the four }\pm1\text{ aligned}}=4,
\]
and by the same alignment $\hat v_6=4$. For a generic other bin, say $k=1$, the
exponential is $e^{-i\pi n/4}$ and the four nonzero samples
$v_0,v_2,v_4,v_6=(+1,-1,+1,-1)$ multiply phases that sum to zero:
\[
  \hat v_1=(1)(1)+(-1)(e^{-i\pi/2})+(1)(e^{-i\pi})+(-1)(e^{-i3\pi/2})
         =1+i-1-i=0,
\]
and likewise $\hat v_0=\hat v_3=\hat v_4=\hat v_5=\hat v_7=0$.

\medskip
\emph{Read off the resonance.} The magnitude spectrum is
\[
  |\hat v_k|=(0,\,0,\,4,\,0,\,0,\,0,\,4,\,0),
\]
a single peak at $k=2$ (and its real-signal mirror at $k=6$). The transform has
located the ring's frequency---two cycles per window---as cleanly as
\cref{fig:tr-fourier} drew it. On a real port signal the ring is not a single pure
tone but a band, and the peak has finite width set by the decay rate (a slow decay,
a high-$Q$ resonance, gives a sharp peak); dividing the output spectrum by the input
pulse's spectrum then yields the S-parameter across the whole band. But the essential
move is exactly the one we just did by hand: a time-domain march, followed by a
Fourier transform, recovers the frequency-domain answer the driven analysis
(\cref{ch:driven}) computes one frequency at a time. One pulsed fold, one transform,
the whole band.
\end{workedexample}

\summarysection
The transient analysis computes the time-domain response of a device to an
excitation pulse $\vJ(t)$: not a snapshot but a \emph{movie}, the whole field
history. It runs the anchor skeleton of \cref{ch:situating} with one corner swapped.
\emph{Assemble} (\S\ref{sec:tr-assemble}): three operators $\Kop,\Cop,\Mop$ on
$\Hcurl$ by three \code{fe\_assemble} folds, plus an implicit ODE integrator and a
seed state---all built \emph{once}, before the march (\cref{ch:time}). \emph{Solve}
(\S\ref{sec:tr-solve}): the \code{fold\_solve} corner, the one genuinely sequential
driver. Where the electrostatic and driven solves are \texttt{map}s---independent
families, parallel over terminals or frequencies---the transient march is a
\texttt{fold}: \code{foldl} of an opaque per-step operator over the timestep
schedule, threading a single field-state carry forward, each step beginning where the
last ended, each step hiding an implicit \code{ksp\_solve} (\cref{ch:krylov}). The
whole feature is the one-line composition
$\texttt{transient}=\code{fold\_solve}\circ\code{fe\_assemble}$. The fold's
schedule-split law makes checkpoint-and-restart exact---the carry is the only thing
crossing a split point. \emph{Reduce} (\S\ref{sec:tr-reduce}): a per-step readout of
fields, port waveforms, and energy, producing the whole trajectory; and, by Fourier
transform of a port signal (\S\ref{sec:tr-fourier}), the broadband bridge to the
frequency-domain S-parameters of \cref{ch:driven}. \textbf{Transient is the one
sequential driver: a fold where each step begins where the last ended---the opposite
of the map drivers.}

\furtherreading
The finite-element time-domain formulation, the second-order semi-discrete system,
and the specific implicit integrators (the generalized-$\alpha$ family and its
numerical damping) are treated in Jin's finite-element electromagnetics
text~\citep{jin2014} in the exact form Palace uses. The time-integration machinery
this chapter composes---the stiffness phenomenon, A-stability, the per-step implicit
solve---is developed in \cref{ch:time}, building on the matrix-computation and
eigenvalue-spread material of Golub and Van Loan~\citep{golubvanloan2013}, who also
cover the dense and sparse linear algebra behind each step's solve. For the physics
of the time-domain response and the energy bookkeeping, Griffiths~\citep{griffiths2017}
is the standard undergraduate source. The map-versus-fold distinction that organizes
this chapter is the functional programmer's daily bread, built from scratch in
\cref{ch:combinators}; the fixed-operator map it is contrasted against is the anchor
of \cref{ch:electrostatics}, and the frequency-domain analysis it bridges to is
\cref{ch:driven}.

\exercisessection
\begin{enumerate}[nosep]
  \item \textbf{Map versus fold, by the picture.} Redraw \cref{fig:tr-foldvmap} from
  memory. Which arrows are present in the fold panel but absent in the map panel, and
  what do those arrows represent physically? State, in one sentence each, why the map
  can be run on three machines at once and why the fold cannot.
  \item \textbf{The one swapped corner.} Using \cref{tab:es-swaps}, describe the
  transient analysis as ``the electrostatic anchor with these stages changed.''
  Exactly which cells move (FE space, solve corner, reduction), and which stay? Which
  single change is the structural heart of the chapter?
  \item \textbf{Two independent costs.} The notation box of \S\ref{sec:tr-cost}
  distinguishes the \emph{carry} cost from the \emph{implicit-solve} cost. Classify
  each of the following as affected by one, the other, both, or neither: (a) doubling
  the number of CPU cores; (b) switching from an implicit to a (hypothetical stable)
  explicit integrator; (c) halving $\Delta t$; (d) refining the mesh so the largest
  eigenvalue grows.
  \item \textbf{One implicit step by hand.} For the two-DoF system of
  \cref{wex:tr-twodof}, take \emph{one} backward-Euler step of size $\Delta t=0.2$
  (rather than $0.1$) from rest. Form $\mat A=\mat I+\Delta t\,\Cop+\Delta t^2\Kop$,
  write the right-hand side, and solve the $2\times2$ system for $\vv_1$. Confirm
  $\mat A$ depends only on $\Delta t$ and the operators, not on the step index---the
  operator-capture-once property.
  \item \textbf{Checkpoint exactness.} Using the schedule-split law \cref{eq:tr-split},
  argue that a transient run checkpointed after every $K$ steps and resumed produces
  \emph{exactly} the same final state as an uninterrupted run, for any $K$. What
  single quantity must the checkpoint save? Why does an electrostatic
  \code{solve\_family} (\cref{ch:electrostatics}) have no analogous ``resume''?
  \item \textbf{The DFT bridge, varied.} Repeat \cref{wex:tr-dft} for the $N=8$ signal
  $v_n=\cos(2\pi\cdot1\cdot n/8)$ (one cycle per window instead of two). Predict which
  bins $|\hat v_k|$ are nonzero \emph{before} computing, then verify by evaluating
  $\hat v_1$ and $\hat v_3$. What does the shifted peak say about the resonant
  frequency the march recorded?
  \item \textbf{Why a pulse?} Explain, in terms of the harmonic rule
  (\cref{ch:reductions-pde}), why a single \emph{narrow} pulse in time lets one
  transient run recover a \emph{wide} band of the frequency response, whereas the
  driven analysis (\cref{ch:driven}) needs one solve per frequency. Give one situation
  where you would nonetheless prefer the driven sweep.
  \item \textbf{The opaque step.} The chapter treats \code{time\_step\_op} as opaque,
  quantifying over it rather than rendering it (\cref{fig:tr-innersolve}). What does
  the fold combinator actually \emph{need} to know about one step for the
  schedule-split law to hold, and what can it safely \emph{not} know? Relate your
  answer to why the same \code{fold\_solve} combinator can also drive the
  adaptive-mesh-refinement loop (\cref{ch:situating}), whose per-step body is an
  entire solve rather than an ODE step.
\end{enumerate}
